
\documentclass[11pt,a4paper]{article}

\usepackage[margin=0.75in]{geometry}
\usepackage{setspace}
\usepackage{enumitem}
\usepackage{booktabs}
\usepackage{tabularx}
\usepackage{xcolor}
\usepackage{hyperref}
\usepackage{titlesec}
\usepackage{lmodern}
\usepackage{parskip}

\definecolor{darkblue}{RGB}{20,45,80}
\definecolor{green}{RGB}{25,120,65}

\hypersetup{
    colorlinks=true,
    linkcolor=green,
    urlcolor=green
}

\titleformat{\section}
  {\large\bfseries\color{darkblue}}
  {\thesection}{0.6em}{}

\titleformat{\subsection}
  {\normalsize\bfseries\color{green}}
  {\thesubsection}{0.6em}{}

\setstretch{1.08}

\begin{document}

\begin{center}
{\LARGE\bfseries\color{darkblue} VendorFlow: A Unified SaaS Platform for Local Grocery \& Pharmacy Digitization}\\[10pt]
{\large Project Proposal for UCS503P}\\[8pt]
\textbf{Submitted to: Dr. Jeelani}\\
Thapar Institute of Engineering and Technology\\
August 16, 2026\\[10pt]
\textbf{Project Team}\\
Pranjal Chitravanshi \quad Kunal Thakur \quad Mannat Saini
\end{center}

\hrule
\vspace{8pt}

\section{Higher-order Goal}

This project aims to strengthen local retail resilience by giving small grocery stores and
pharmacies an affordable path to digital commerce. While the broader vision is economic
empowerment of neighborhood businesses, the immediate engineering objective is to translate
reliable order fulfillment and vendor coordination into a measurable, deployable software
solution.

\section{Time-to-value}

\begin{itemize}[leftmargin=1.5em]
    \item We will deliver meaningful increments quickly by prototyping the core ordering and
    vendor-routing workflow and validating it with a small set of pilot vendors within a short
    iteration window.
    \item We will prioritize fast feedback over perfection by using weekly iterations and
    targeting CI-backed automation early to reduce release friction.
    \item Success is defined by what can be delivered and measured in the first iteration, then
    improved based on vendor and customer feedback.
\end{itemize}

\section{Problem Statement}

Small grocery stores and pharmacies in most Indian cities and towns rely on phone calls,
WhatsApp, spreadsheets, or paper records to manage orders and inventory. Common pain points
include:

\begin{itemize}[leftmargin=1.5em]
    \item \textbf{No online presence:} Most vendors cannot afford a custom website, app, or
    delivery infrastructure of their own.
    \item \textbf{Manual, error-prone operations:} Inventory and order tracking done by hand
    leads to stockouts, missed orders, and delivery mix-ups.
    \item \textbf{Fragmented discovery:} Customers have no single place to find nearby stores,
    compare stock, or track deliveries.
    \item \textbf{Regulatory friction for pharmacies:} Prescription-backed orders need a
    verification step that ad-hoc channels such as calls and chat apps do not support reliably.
\end{itemize}

Consumer demand for online ordering continues to grow, but the smallest vendors are the
ones least able to build this infrastructure themselves. A shared platform lets many small
businesses digitize at once without each bearing the full engineering cost.

\section{Proposed Solution}

\subsection{Overview}

VendorFlow is proposed as a full-stack SaaS platform with the following components:

\begin{itemize}[leftmargin=1.5em]
    \item \textbf{Customer portal:} Discover nearby grocery/pharmacy vendors, browse products,
    upload prescriptions where required, place orders, and track deliveries.
    \item \textbf{Vendor portal:} Vendor onboarding and approval workflow, inventory management,
    order management, and reliability/analytics dashboard.
    \item \textbf{Delivery partner portal:} Accept assigned deliveries and update status in real time.
    \item \textbf{Admin portal:} Vendor approval, prescription verification oversight, platform
    analytics, and audit logs.
    \item \textbf{Smart vendor routing:} Route an order to the best-fit vendor using distance,
    live stock, and a computed reliability score.
    \item \textbf{Credit system:} An in-platform credit/wallet mechanism for customers and vendors,
    supporting promotional credits, refunds, and settlement between platform and vendors.
\end{itemize}

\subsection{Core Workflow}

\begin{enumerate}[leftmargin=1.8em]
    \item The customer browses nearby vendors and adds items to the cart, uploading a prescription
    for pharmacy items where required.
    \item The system applies smart vendor routing using distance, stock, and reliability score.
    For pharmacy orders, the prescription is queued for manual verification.
    \item The order is confirmed, a delivery partner is assigned, and status updates flow through
    the lifecycle: \textbf{Placed $\rightarrow$ Confirmed $\rightarrow$ Out for Delivery $\rightarrow$ Delivered}.
    \item The customer and vendor receive real-time notifications at each stage. Credit/wallet
    balances are updated on completion, refund, or promotional events.
\end{enumerate}

\subsection{Operational Constraints}

\begin{itemize}[leftmargin=1.5em]
    \item Keep customer- and vendor-facing interfaces usable on low-to-mid range devices through
    responsive web design.
    \item Keep order-matching and routing logic heuristic and explainable rather than research-grade,
    focusing on correctness, latency, and auditability.
    \item Design the backend for high throughput because order placement, routing, and notification
    delivery are on the platform's critical path.
\end{itemize}

\section{Solution Approach}

This is an engineering course project, so the focus is on system architecture, performance,
and correctness of the transactional workflow rather than novel algorithm research.

\subsection{Performance-focused Backend}

\begin{itemize}[leftmargin=1.5em]
    \item \textbf{C++:} High-throughput backend services including order intake, routing engine,
    and inventory locking for predictable low-latency performance under load.
    \item \textbf{Python:} Auxiliary services such as analytics jobs, reliability-score
    recomputation, and admin/reporting tooling.
    \item \textbf{JavaScript:} Customer, vendor, delivery, and admin front-ends, together with
    lightweight orchestration/API-gateway glue where appropriate.
    \item \textbf{Vendor routing:} A simple weighted-score heuristic based on distance, stock
    availability, and reliability score, without claiming state-of-the-art optimization.
\end{itemize}

\subsection{System Architecture}

The proposed system follows a multi-tier architecture:

\begin{itemize}[leftmargin=1.5em]
    \item \textbf{Frontend:} Responsive web applications for customer, vendor, delivery partner,
    and admin portals, built in JavaScript.
    \item \textbf{High-performance backend:} C++ REST/WebSocket services for authentication,
    order lifecycle, vendor routing, inventory locking, and the credit/wallet ledger.
    \item \textbf{Auxiliary services:} Python background jobs for reliability-score updates,
    routing-timeout handling, notification queue processing, and analytics pipelines.
    \item \textbf{Cache/queue layer:} Redis for caching, rate limiting, session data, and
    background job/notification queues.
    \item \textbf{Primary database:} PostgreSQL for users, vendors, products, inventory, orders,
    credit/wallet ledger, reviews, and audit logs.
    \item \textbf{Real-time layer:} Socket.IO or an equivalent WebSocket layer for real-time
    order and delivery status notifications.
\end{itemize}

\subsection{CI/CD}

CI/CD will be used to:

\begin{itemize}[leftmargin=1.5em]
    \item automatically build and run C++ unit tests, Python service tests, and frontend tests
    on every change;
    \item produce repeatable deployment artifacts to staging for each service; and
    \item enable safe and frequent releases of incremental features across the portals.
\end{itemize}

\section{Evaluation Criteria}

Success will be evaluated using measurable metrics tied to the core order-fulfillment workflow.

\subsection{Primary Evaluation Metric}

\textbf{Order-to-confirmation latency (OCL)} is the time between order placement and vendor
confirmation after routing.

\begin{itemize}[leftmargin=1.5em]
    \item \textbf{Target:} Median OCL $\leq$ 60 seconds during the pilot window.
    \item \textbf{Attribution:} Measured using backend timestamps from the order-placed event
    to the vendor-confirmed event.
\end{itemize}

\subsection{Secondary Evaluation Metrics}

\begin{itemize}[leftmargin=1.5em]
    \item Delivery success rate: percentage of orders marked Delivered without escalation or cancellation.
    \item Vendor reliability accuracy: correlation between the computed reliability score and
    actual on-time fulfillment during the pilot.
    \item Prescription verification turnaround: median time from upload to admin-verified status
    for pharmacy orders.
    \item System reliability: availability of login, ordering, and routing during business hours,
    with an example pilot target of at least 99\% uptime.
    \item Backend throughput: sustained requests per second served by the C++ order and routing
    services under load-test conditions.
\end{itemize}

\subsection{Pilot Validation Plan}

The pilot will recruit approximately 10--20 customers and 3--5 vendor accounts, with a mix of
grocery and pharmacy vendors. Where real vendors are unavailable, simulated vendors may be used.
Metrics will be compared before and after implementation using short interviews or existing
process observations to estimate baseline manual-process timings.

\section{Scalability}

The system should scale both theoretically and practically.

\subsection{Theoretical Scalability}

The system should support growth in vendors, orders, and concurrent portal usage by:

\begin{itemize}[leftmargin=1.5em]
    \item decoupling the C++ backend, Python auxiliary services, and frontends;
    \item enabling pagination, indexing, and read replicas for common PostgreSQL queries;
    \item using Redis for caching and rate limiting to shield the primary database; and
    \item designing backend APIs to be stateless where possible.
\end{itemize}

\subsection{Operational Deployability}

The system should run on common infrastructure using:

\begin{itemize}[leftmargin=1.5em]
    \item managed PostgreSQL and Redis instances;
    \item containerized deployment for the C++ backend, Python services, and frontends; and
    \item a CI/CD pipeline for staging and production across all services.
\end{itemize}

\section{Technology and Engine Availability}

The project will use mature technologies and services rather than inventing new infrastructure:

\begin{itemize}[leftmargin=1.5em]
    \item High-performance C++ web/service framework for the core backend.
    \item Python framework for auxiliary/background services and analytics.
    \item JavaScript framework for the four frontend portals.
    \item PostgreSQL as the managed relational database.
    \item Redis for caching, rate limiting, and job/notification queues.
    \item JWT-based authentication with role-based access control.
    \item Object storage for prescription images and product photos.
    \item Unit and integration testing frameworks across C++, Python, and JavaScript.
    \item CI runner and staging deployment target.
\end{itemize}

\section{Project Scope and Deliverables}

\subsection{Initial Deliverable}

\begin{itemize}[leftmargin=1.5em]
    \item Customer, Vendor, Delivery Partner, and Admin portals with JWT authentication and
    role-based access control.
    \item Vendor onboarding and approval workflow with product and inventory management.
    \item Smart vendor routing based on distance, stock, and reliability implemented in the C++ backend.
    \item Shopping cart, checkout, and order lifecycle management.
    \item Prescription upload with manual admin verification.
    \item Delivery assignment and status-based tracking.
    \item Basic logging and timestamps for OCL measurement.
    \item CI with builds, C++/Python backend unit tests, and linting.
    \item CD to staging with a single pipeline job.
\end{itemize}

\subsection{Subsequent Deliverables}

\begin{itemize}[leftmargin=1.5em]
    \item Credit/wallet system for customers and vendors, including promotional credits,
    refunds, and settlement.
    \item Real-time notifications through Socket.IO, together with email and in-app notifications.
    \item Reviews and ratings feeding into the vendor reliability score.
    \item Vendor and admin analytics dashboards.
    \item Redis-backed caching and rate limiting hardened for scale.
    \item Background jobs for notification queueing, routing timeouts, and reliability-score recomputation.
    \item Audit logs and activity history across vendor and admin actions.
    \item Improved search/filtering and indexed queries for scale readiness.
\end{itemize}

\section{Risks and Mitigations}

\begin{itemize}[leftmargin=1.5em]
    \item \textbf{Data privacy and consent:} Minimize collected data, especially prescription images;
    store only what is needed and restrict access to verified administrators.
    \item \textbf{Vendor adoption:} Keep the vendor portal simple, with minimal steps for onboarding,
    inventory updates, and order confirmation.
    \item \textbf{Backend complexity:} Mitigate C++ complexity through strong test coverage,
    clear service boundaries, and CI-enforced builds before merge.
    \item \textbf{Evaluation measurement ambiguity:} Instrument timestamps and define clear
    order-lifecycle states before development begins.
    \item \textbf{Operational issues during the pilot:} Use staging early and ensure CI/CD produces
    repeatable deployments for all services.
\end{itemize}

\section{Timeline}

\begin{center}
\renewcommand{\arraystretch}{1.25}
\begin{tabularx}{\textwidth}{>{\bfseries}p{0.19\textwidth} X X}
\toprule
Phase & Major Tasks & Expected Deliverable \\
\midrule
Weeks 1--2 & Requirements, architecture, database design, and UI planning &
System requirements and architecture design \\
Weeks 3--4 & Authentication, role-based access, vendor onboarding, and inventory &
Initial customer/vendor/admin foundation \\
Weeks 5--6 & Cart, checkout, order lifecycle, and smart vendor routing &
Core ordering and routing workflow \\
Weeks 7--8 & Prescription upload, manual verification, and delivery assignment &
Pharmacy and delivery workflow \\
Weeks 9--10 & Real-time status, Redis queues/caching, wallet/credit functionality &
Integrated transactional workflow \\
Weeks 11--12 & Analytics, reviews, testing, CI/CD, load testing, and refinement &
Validated prototype and staging deployment \\
\bottomrule
\end{tabularx}
\end{center}

\section{Conclusion}

VendorFlow is a deployable SaaS platform intended to help small grocery stores and pharmacies
digitize inventory, ordering, delivery, and vendor management. The proposed architecture uses
a high-performance C++ backend, Python auxiliary services, PostgreSQL and Redis for storage
and caching, and JavaScript frontends across four portals.

The project emphasizes time-to-value through rapid iterations, CI/CD for frequent and safe
releases, and measurable evaluation criteria, particularly order-to-confirmation latency.
The focus remains on engineering workflow, backend performance, correctness, and scalability
readiness rather than research-driven novelty.

\end{document}
